\section*{Chapter \ref{ch:fm:case-studies-i} --- Case Studies I}

\begin{solution}{exo:fm:case-studies-i:1}
$2.3/(1-0.44)=\$4.11$ billion at the end of July. On 28 September the \$3.6 billion was 90\% of the net asset value, so the existing investors held
$3.6/0.9-3.6=\$0.40$ billion: a return of $0.40/2.3-1=-82.6\%$ since the end of August.
\end{solution}

\begin{solution}{exo:fm:case-studies-i:2}
See \cref{def:fm:case-studies-i:liquidity}. In 1998, funding: lenders and counterparties could call margin faster than the spreads recovered. In
2006, market: the fund held most of the open interest in some contracts and had no one to sell to. In August 2007, market liquidity through
crowding, turned into funding pressure by leverage and risk limits.
\end{solution}

\begin{solution}{exo:fm:case-studies-i:3}
$100\,000/(0.1\times25\,000)=40$ days.
\end{solution}

\begin{solution}{exo:fm:case-studies-i:4}
$1/28=3.6\%$; the balance sheet was about $28\times4.67=\$131$ billion. Net of hedges, the spread positions moved far less than their gross value on
a normal day, which is why the leverage looked acceptable.
\end{solution}

\begin{solution}{exo:fm:case-studies-i:5}
See the proof of \cref{prop:fm:case-studies-i:days}. At 75\% of the open interest the volume available to the position is small against its size,
so any participation low enough to limit impact means weeks of exposure, and the fund's own sales are most of the market's selling.
\end{solution}

\begin{solution}{exo:fm:case-studies-i:6}
Through day 5 the holder keeps its whole book while the seller's shrinks, so the holder marks more of each day's move. After the selling stops the
temporary impact decays: the holder recovers on its whole book, the seller only on the half it kept, having sold the rest at the moved prices.
\end{solution}

\begin{solution}{exo:fm:case-studies-i:7}
At an overlap of 0.5 (0.52 measured) the seller is down 9.8\% at day 20 and the holder 4.3\%; with unrelated books the seller is down 9.7\% and the
holder 0.1\% up. The seller's loss is its own impact; the holder's is the overlap.
\end{solution}

\begin{solution}{exo:fm:case-studies-i:8}
A one-day value at risk assumes the position can be exited within a day near the model's prices. A position of weeks to liquidate carries weeks of
price risk, and its own sales move the price; the natural-gas fund's positions were most of the open interest in some contracts.
\end{solution}

\begin{solution}{pb:fm:case-studies-i:1}
\begin{enumerate}
\item The public record and the numbers reconstructed from it; it avoids villains, hindsight and motives the record does not state.
\item See \cref{def:fm:case-studies-i:liquidity}.
\item Losses raise margin calls, calls force sales, sales move prices, and moved prices are losses for everyone holding the same positions
(\cref{fig:fm:case-studies-i:spiral}).
\item See \cref{def:fm:case-studies-i:days}.
\item See \cref{dat:fm:case-studies-i:record}: \$4.67 billion of capital at the end of 1997 at 28 to 1, \$2.3 billion at the end of August after a
44\% loss, and \$3.6 billion from 14 firms on 28 September, 90\% of the net asset value.
\item It judged that a rapid liquidation of the fund's positions, and of related positions of others, might threaten already unsettled markets.
\item Up to 100\,000 contracts in a month, about 40\% of the winter 2006--07 NYMEX contracts and 75\% of November's; more than \$2 billion lost from
late August to mid-September, and the natural gas portfolio liquidated.
\item Overlapping trades, higher volatility and correlations; risk and leverage reduced; \$3 billion of equity into a quantitative fund worth about
\$3.6 billion.
\item $-12.1\%$ January to July, $-44\%$ in August, $-82.6\%$ to 28 September.
\item Actual \$4.67, 4.11, 2.30 and 0.40 billion; at half the leverage \$4.67, 4.39, 3.42 and 2.01 billion.
\item The same asset returns, when part of September's fall was trading against a fund in distress; no credit for fewer margin calls; and nothing taken
from the years before, which half the leverage would have halved.
\item 20 days and 40 days, at the illustrative 25\,000 contracts a day.
\item See \cref{prop:fm:case-studies-i:days}.
\item Seller $-5.4\%$, $-13.6\%$, $-9.7\%$; holder $-5.3\%$, $-15.1\%$, $-7.9\%$; the holder keeps more of the book through the selling and so more of the
recovery after it.
\item Small-looking risk at a size and leverage that made the exit a market event, and dependence on lenders, the other side and other funds.
\item A pairing of leverage with funding term, limits on share of open interest and \omterm{def:fm:case-studies-i:days}{days to liquidate}, and a crowded-exit stress.
\item Estimate the overlap of the book with public holdings or crowding measures (Book 7), then run an unwind of a large overlapping seller and
charge the book for the moves.
\item Leverage sets how far losses reach the capital; the funding term sets how long the firm can wait for them to reverse.
\item \$2.01 billion instead of \$0.40 billion at half the leverage; 20 days at 20\% of the illustrative volume (40 at 10\%); the seller down 9.7\% at day
20 against 7.9\% for the holder.
\item Pair every leverage limit with a minimum funding term and cap each position's \omterm{def:fm:case-studies-i:days}{days to liquidate} and share of open interest. Stress the book for a
crowded exit, and decide in advance whether it would sell or hold.
\end{enumerate}
\end{solution}

\begin{solution}{iq:fm:case-studies-i:1}
Funding: a lender raises haircuts and the firm cannot meet the call. Market: a position is so large against volume that selling it moves the price.
\iqlookfor{both definitions and how they feed each other.}
\end{solution}

\begin{solution}{iq:fm:case-studies-i:2}
$1/25=4\%$. Gross assets overstate the risk of hedged spread positions; the relevant measure is the net risk under stress, when spreads move together.
\iqlookfor{gross against net.}
\end{solution}

\begin{solution}{iq:fm:case-studies-i:3}
By the \omterm{def:fm:case-studies-i:days}{days to liquidate} and the impact of an exit, not by a one-day risk measure; limit the share of open interest and the \omterm{def:fm:case-studies-i:days}{days to liquidate}.
\iqlookfor{exit cost over one-day risk.}
\end{solution}

\begin{solution}{iq:fm:case-studies-i:4}
It depends on funding and limits: if the loss is a crowded exit and the firm can hold through it, holding recovers more; if a call or limit will
force the cut later, cutting early costs less. Check the overlap and the funding horizon first.
\iqlookfor{the decision depends on funding, not only on the signal.}
\end{solution}

\begin{solution}{iq:fm:case-studies-i:5}
A forced liquidation would move prices against their own positions in the same trades and against their exposures to the fund.
\iqlookfor{the creditors' own books.}
\end{solution}

\begin{solution}{iq:fm:case-studies-i:6}
The book's positions, their volumes and volatilities, an impact model, and the size and overlap of other holders' books, which it must estimate from
public holdings, crowding measures or past episodes.
\iqlookfor{the unobserved overlap.}
\end{solution}
